\section{Conclusion}

Ce laboratoire a permis d'étudier un circuit magnétique double~E avec entrefer réglable, en confrontant un modèle analytique aux mesures expérimentales.

\paragraph{Partie théorique.}
Le modèle développé prend en compte les \textbf{trois entrefers} de la structure double~E (colonne centrale + deux colonnes latérales), ce qui conduit à la réluctance équivalente correcte :
\[
\mathcal{R}_{eq}=\frac{3(x+\lambda)}{2\mu_0 S_g},
\qquad
B_\mathrm{lat}=\frac{\mu_0 N_1 I}{3(x+\lambda)}.
\]
Les résultats théoriques principaux sont :
\begin{itemize}
\item \(R_1(20^\circ\mathrm{C})\approx\SI{6.65}{\ohm}\),
      \(R_1(100^\circ\mathrm{C})\approx\SI{8.74}{\ohm}\),
      \(R_1(150^\circ\mathrm{C})\approx\SI{10.05}{\ohm}\) ;
\item \(I_\mathrm{ref}\approx\SI{0.50}{A}\) pour \(J=\SI{4}{A/mm^2}\) ;
\item \(L_{12}(0\,\mathrm{mm})\approx\SI{34.6}{mH}\), \(L_{12}(\SI{2}{mm})\approx\SI{5.8}{mH}\) ;
\item \(f_c\approx\SI{1.59}{Hz}\) pour l'intégrateur (\(R_{10}C=\SI{0.1}{s}\)).
\end{itemize}

\paragraph{Partie pratique — résultats obtenus.}
\begin{itemize}
\item \textbf{Résistance} : \(R_1=\SI{6.164}{\ohm}\) (4 fils), écart de \(-7.3\,\%\) par rapport à la théorie, expliqué par les tolérances de fabrication du fil.
\item \textbf{Induction \(B=f(I)\)} : comportement linéaire sur \(0\)–\(\SI{0.6}{A}\), accord théorie/mesure à \(\pm 5\,\%\) (\(x=\SI{2}{mm}\)) et \(\pm 9\,\%\) (\(x=\SI{3}{mm}\)), ce qui valide le modèle à 3 entrefers avec \(\lambda=\SI{0.4}{mm}\).
\item \textbf{Intégrateur} : le diagramme de Bode confirme \(f_c\approx\SI{1.59}{Hz}\) et le régime quasi-intégrateur à \(\SI{35}{Hz}\).
\item \textbf{Inductance mutuelle \(L_{12}\)} : les valeurs mesurées (\SI{139}{mH} à \(x=0\), \SI{14.5}{mH} à \(x=\SI{2}{mm}\)) sont significativement supérieures aux prédictions théoriques (+302\,\% et +150\,\% respectivement). Ces écarts s'expliquent par la perméabilité réelle du matériau bien supérieure à 400 et par les effets de franges qui augmentent la section d'entrefer effective.
\end{itemize}

\paragraph{Discussion.}
Les écarts résiduels entre théorie et mesure s'expliquent par :
\begin{itemize}
\item les \textbf{effets de franges} (flux de fuite autour des entrefers), qui augmentent la perméance effective et augmentent \(B\) par rapport au modèle idéal ;
\item la \textbf{perméabilité réelle} du matériau (\(\mu_r\) non-constante, dépendant du point de fonctionnement) ;
\item les \textbf{tolérances mécaniques} sur l'entrefer \(x\) et les dimensions du circuit.
\end{itemize}
Pour les petits entrefers, une saturation partielle pourrait apparaître à courant élevé, ce que les mesures commencent à révéler par une légère sous-linéarité.

\paragraph{Remarque sur le tableau AC.}
Le tableau des mesures AC (Table~\ref{tab:AC_ref}) reste à compléter lors de la séance de mesures. Les valeurs de \(L_{12}\) ont été estimées à partir du diagramme \(x\)-\(y\) de l'oscilloscope selon \(L_{12}=(\Delta V_\mathrm{out}/\Delta I_1)\cdot R_9 C\).
